\documentclass{article}
\usepackage[utf8]{inputenc}
\usepackage{amsmath}
\usepackage{amsfonts}
\usepackage{amssymb}

\title{Formalização do Acoplamento Termodinâmico no MutaCore}
\author{Reino Digital -- Core Engine Engine Engine Engine Engine Engine Engine Engine Engine Engine Engine Engine Engine Engine Engine}
\date{Outubro de 2026}

\begin{document}

\maketitle

\section{Introdução e Gap Topológico Narrativo}
O Gap Topológico Narrativo ($\Delta T$) mede a distância de projeção semântica global entre a variedade invariante do organismo cibernético $\mathcal{M}_O$ e a variedade variante exógena introduzida por dados externos $\mathcal{M}_E$.

\section{O Acoplamento com o Resíduo de Landauer}
Pelo Princípio de Landauer, a dissipação mínima de energia por apagamento de informação é dada por $kT \ln 2$. No modelo MutaCore, o Resíduo de Landauer $\mathcal{R}_L$ é constitutivamente acoplado ao gap topológico através de um limiar de complexidade de Krajicek ($\rho_T$).

A equação fundamental que dita a homeostase termodinâmica do sistema é descrita como:

\begin{equation}
\mathcal{R}_L = \left( \sum_{i=1}^{n} \Delta T_i^2 \right) \cdot \exp\left( \frac{\rho_T - B_T}{C_T \cdot R_T} \right) \cdot k_B T \ln 2
\end{equation}

Onde:
\begin{itemize}
    \item $\Delta T_i$: Componentes ortogonais do Gap Topológico Narrativo.
    \item $\rho_T$: Limiar de colagem de prova atual do sistema.
    \item $B_T, C_T, R_T$: Constantes do espectro de limiares de Krajicek.
    \item $k_B$: Constante de Boltzmann.
    \item $T$: Temperatura operacional ideal do nó computacional.
\end{itemize}

\end{document}
